\documentclass[11pt,a4paper]{article}

\usepackage[utf8]{inputenc}
\usepackage[T1]{fontenc}
\usepackage{lmodern}
\usepackage[margin=1in]{geometry}
\usepackage{booktabs}
\usepackage{tabularx}
\usepackage{xcolor}
\usepackage{enumitem}
\usepackage{hyperref}
\usepackage{microtype}

\definecolor{primaryblue}{RGB}{24,76,120}
\definecolor{neutraldark}{RGB}{35,39,42}

\hypersetup{
    colorlinks=true,
    linkcolor=primaryblue,
    urlcolor=primaryblue,
    pdftitle={BTP Project Progress Update},
    pdfauthor={G Santhosh Reddy, G Vamsi Krishna, S Sri Nithish}
}

\setlist[itemize]{leftmargin=1.5em,itemsep=3pt,topsep=4pt}
\setlist[enumerate]{leftmargin=1.7em,itemsep=3pt,topsep=4pt}

\begin{document}

\begin{center}
    {\Large\bfseries\color{primaryblue} BTP Phase II: Project Progress Update}\\[5pt]
    {\large Privacy Red-Teaming for Unstructured Data in Multi-Agent Systems}\\[8pt]
    {\normalsize Indian Institute of Technology Tirupati \textbar\
    Department of Computer Science and Engineering}\\[4pt]
    {\normalsize 2 November 2026}
\end{center}

\vspace{8pt}
\noindent\textbf{Project team:}
G Santhosh Reddy (CS23B022), G Vamsi Krishna (CS23B023), and
S Sri Nithish (CS23B052).

\section{Executive Summary}

We have implemented a working software prototype for studying privacy risks
when multi-agent systems process unstructured documents. The prototype includes
document ingestion, PII detection and redaction, persistent chunk retrieval,
sequential and hierarchical agent workflows, and an adversarial benchmark
harness. The current automated test suite passes \textbf{74 tests}.

The benchmark generator and evaluation pipeline are implemented, but a live
evaluation using Ollama models has not yet been completed. Accordingly, this
update reports implementation and test status only; it does not claim empirical
LLM performance or measured leakage reduction.

\section{Work Completed}

\begin{description}[leftmargin=0pt,labelindent=0pt,style=nextline]
    \item[\textcolor{primaryblue}{Document ingestion and extraction}]
    A FastAPI application accepts documents and supports text extraction from
    PDF, DOCX, and TXT files. Uploaded documents can be stored locally and
    passed into the agent workflow.

    \item[\textcolor{primaryblue}{PII privacy controls}]
    The privacy guardrail combines Presidio-based analysis, when available,
    with regular-expression detection. It supports redaction, risk scoring,
    and before/after telemetry. A separate egress gateway inspects outbound
    payloads and records its decision and receipt, distinguishing attempted
    PII exposure from PII allowed past that gateway.

    \item[\textcolor{primaryblue}{Persistent retrieval and data-flow auditing}]
    ChromaDB-backed vector storage supports document chunking and semantic
    retrieval. A chunk-flow auditor records which workflow roles handle
    chunks and checks whether chunks reaching the analyst passed through the
    privacy-redaction step.

    \item[\textcolor{primaryblue}{Multi-agent workflows}]
    The application supports both a sequential CrewAI workflow and a
    hierarchical workflow with a supervisor audit step. Guardrail-enabled
    execution and a sequential-versus-hierarchical comparison endpoint are
    also available.

    \item[\textcolor{primaryblue}{Adversarial benchmark infrastructure}]
    A seeded fuzzer generates synthetic documents in PDF, DOCX, and TXT
    formats across four example domains: HR, finance, medical, and legal.
    It covers five indirect-injection threat classes: context hijacking,
    selective exfiltration, delimiter smuggling, role confusion, and encoded
    payloads. The evaluator is configured for three Ollama models
    (\texttt{llama3.2}, \texttt{mistral}, and \texttt{phi3}) and three
    architecture configurations (baseline, guardrail-enabled, and
    defense-in-depth).

    \item[\textcolor{primaryblue}{Benchmark telemetry and access}]
    The evaluation pipeline records benign-task success, injection success,
    attempted PII exposure, and guardrail redaction. Gateway receipts separately
    capture whether an attempted payload was allowed and how much PII was
    transmitted. API endpoints and a command-line runner are available to
    generate corpora, execute runs, and save result summaries.
\end{description}

\section{Verification and Current Status}

On 2 November 2026, the project test suite was run with
\texttt{uv run pytest -q}: \textbf{74 passed}. The tests cover core workflows,
document extraction, privacy controls, storage, and benchmark generation and
evaluation logic. The benchmark tests use deterministic offline simulation for
repeatable checks; simulated outputs are not evidence of live model behavior.

\begin{center}
\small
\begin{tabularx}{\textwidth}{@{}l X@{}}
\toprule
\textbf{Area} & \textbf{Status} \\
\midrule
Document processing and agent workflows & Implemented and covered by automated tests \\
PII guardrail, egress inspection, and retrieval auditing & Implemented and covered by automated tests \\
Adversarial corpus generation and evaluation tooling & Implemented; live model validation pending \\
Live Ollama benchmark results & Pending; no empirical results are reported here \\
Expanded repeated evaluation and final analysis & Next work \\
\bottomrule
\end{tabularx}
\end{center}

\section{Next Steps}

\begin{enumerate}
    \item Confirm that the configured Ollama models are available and run a
    live pilot across the three models and three architecture configurations.
    \item Review the resulting run records, including attempted exposure,
    gateway decisions, transmitted PII, and guardrail redaction, before
    drawing conclusions.
    \item Execute the full threat-class matrix with repeated samples, then
    analyze the outcomes and relevant performance trade-offs.
    \item Incorporate validated results into the final project report and
    demonstration.
\end{enumerate}

\vfill
\noindent\textit{\color{neutraldark}This is a separate progress update;
no edits to the existing Stage 3 report were made for this update.}

\end{document}
